% =====================================================================
\section{What device experiments show}
\label{sec:devices}

\pgfplotsset{timelineaxis/.style={ytick={1,...,16}, ymin=0.4, ymax=16.6, yticklabels={{Google Sycamore RCS, 53 q (2019)}, {Xanadu Borealis GBS (2022)}, {USTC Jiuzhang 3.0 GBS (2023)}, {IBM doped Clifford, 70 q (v1, 2026)}, {USTC Jiuzhang GBS, 100 modes (2020)}, {USTC Jiuzhang 2.0 GBS (2021)}, {IBM Eagle kicked Ising, 127 q (2023)}, {D-Wave annealing dynamics (v1, 2024)}, {D-Wave QAC optimization (v1, 2024)}, {Harvard Rydberg MIS, 289 atoms (2022)}, {Quantinuum H2 magnetism, 56 q (2025)}, {Google OTOC(2), 65 q (2025)}, {IBM--Algorithmiq OLE, 56 q (2025)}, {USTC ultracold atoms, 64 sites (v2, 2026)}, {IBM Heron Floquet, 74 q (v1, 2026)}, {IBM doped Clifford, 64 q (v2, 2026)}}}}

\begin{figure}[htbp]
\centering
\begin{tikzpicture}
\begin{axis}[
  xbar, width=0.56\textwidth, height=11.3cm,
  xmin=0, xmax=38, y dir=reverse, timelineaxis,
  yticklabel style={font=\footnotesize},
  xlabel={\footnotesize Months from first public claim (days / 30.44); open claims to 11 Nov 2026},
  xticklabel style={font=\footnotesize},
  bar width=7pt, bar shift=0pt,
  xmajorgrids, grid style={black!12},
  axis x line*=bottom, axis y line*=left, clip=false,
  legend style={font=\scriptsize, at={(0.5,1.02)}, anchor=south, draw=black!25, fill=white, legend columns=3}, legend cell align=left]
\addplot[fill=qcN, draw=black!75, line width=0.3pt, nodes near coords={\pgfmathprintnumber[fixed,precision=1]{\pgfplotspointmeta}\,mo, matched}, point meta=x, every node near coord/.append style={font=\scriptsize, anchor=west}] table[x=x,y=y] {figdata/timeline_validation.dat};
\addlegendentry{validation benchmark matched}
\addplot[fill=white, draw=black!75, line width=0.3pt, postaction={pattern=north east lines, pattern color=black!70}, nodes near coords={\pgfmathprintnumber[fixed,precision=1]{\pgfplotspointmeta}\,mo, challenged}, point meta=x, every node near coord/.append style={font=\scriptsize, anchor=west}] table[x=x,y=y] {figdata/timeline_challenged.dat};
\addlegendentry{validation statistic challenged}
\addplot[fill=black!85, draw=black!85, line width=0.3pt, nodes near coords={\pgfmathprintnumber[fixed,precision=1]{\pgfplotspointmeta}\,mo, reproduced}, point meta=x, every node near coord/.append style={font=\scriptsize, anchor=west}] table[x=x,y=y] {figdata/timeline_observables.dat};
\addlegendentry{observables reproduced}
\addplot[fill=qcNl, draw=black!75, line width=0.3pt, postaction={pattern=dots, pattern color=black!80}, nodes near coords={$\approx$\pgfmathprintnumber[fixed,precision=1]{\pgfplotspointmeta}\,mo, partly}, point meta=x, every node near coord/.append style={font=\scriptsize, anchor=west}] table[x=x,y=y] {figdata/timeline_partial.dat};
\addlegendentry{observables partly reproduced}
\addplot[fill=qcUl, draw=qcUd, line width=0.3pt, postaction={pattern=crosshatch, pattern color=qcU}, nodes near coords={\pgfmathprintnumber[fixed,precision=1]{\pgfplotspointmeta}\,mo, contested}, point meta=x, every node near coord/.append style={font=\scriptsize, anchor=west}] table[x=x,y=y] {figdata/timeline_contested.dat};
\addlegendentry{contested}
\addplot[fill=qcF, draw=qcQ, line width=0.3pt] table[x=x,y=y] {figdata/timeline_open.dat};
\addlegendentry{open (right-censored)}
\addplot[quiver={u=3.2,v=0}, -{Latex[length=2.2mm]}, qcF, line width=1.2pt, forget plot] table[x=x,y=y] {figdata/timeline_open.dat};
\addplot[only marks, mark=none, forget plot, nodes near coords={\textbf{$\ge$\pgfmathprintnumber[fixed,precision=1]{\pgfplotspointmeta}\,mo, open}}, point meta=x, every node near coord/.append style={font=\scriptsize, anchor=west, qcF, xshift=22pt}] table[x=x,y=y] {figdata/timeline_open.dat};
\end{axis}
\end{tikzpicture}
\caption[Time from quantum claim to classical response]{Time from first public claim to first classical response for the \nTimeline{} device claims that meet the eligibility rule: \emph{a claim that a quantum device performed a computation faster or at larger scale than the best stated classical method}. The kind of response (Section~\ref{sec:match}) is printed in words at the end of each bar and repeated in the fill: a plain mid-tone bar for a matched validation benchmark, diagonal hatching for a challenged validation statistic, a plain dark bar for reproduced observables, dots for a partial reproduction, cross-hatching for a contested response, and a bar ending in an arrow for an open claim; arrows mark claims still open on 11 November 2026, from which no rate is inferred. Each row is one instance and version. A short bar records a fast response to a specified test; an open bar records the absence of a qualifying response, not evidence of impossibility. For Jiuzhang~2.0 the bar is the 2021 challenge, and the 2023 match is described in the text. A response to the 70-qubit doped Clifford circuit does not stop the clock for its 64-qubit successor, and the ultracold-atom clock starts at the April 2026 version that first reported 64 sites. The operator-Loschmidt-echo row is plotted from 30 November 2025, a lower bound on its age. Later random-circuit generations and the newest boson-sampling generation are not entered. Dates and sources are in \texttt{entries.csv}.}
\label{fig:timeline}
\end{figure}


Fig.~\ref{fig:timeline} places \nTimeline{} device claims on a common clock. Of these, \wTLvalidation{} were answered by matching the statistic the experiment used to validate itself, \wTLchallenged{} were challenged on such a statistic by classical mock-up samplers, \wTLobservables{} by reproducing its reported observables and \wTLpartial{} in part; \wTLcontested{} are contested and \wTLopen{} have no published classical match as of 11 November 2026. The answered claims took between 12 days and 24 months. The open ones are between \eld{E25} and \eld{E05} days old. The figure is a record of what happened to these claims. It is not a survival curve, and the open bars predict nothing. Each case below is described by the same five questions: what scientific output is at stake, what each side computed, what kind of evidence that is, what remains unresolved, and what result would settle it.

\subsection{Sampling benchmarks}
\textbf{Question.} Can a device draw samples from a probability distribution that no classical computer can sample at comparable quality? These experiments test capability and have no application output.

\textbf{What was done.} Google's Sycamore sampled random 53-qubit circuits in 2019~\cite{arute2019}, and USTC's Jiuzhang sampled photons from a Gaussian boson network from 2020~\cite{zhong2020}; larger random-circuit experiments followed at Google, USTC and Quantinuum. In July 2026 IBM sampled a 70-qubit circuit of Clifford gates doped with 468 T gates, using error-detecting spacetime codes to certify fidelity~\cite{martiel2026}. USTC sampled a driven Bose--Hubbard system of up to 64 lattice sites and 20 atoms in a quantum gas microscope and reported a sampling rate three orders of magnitude above a supercomputer~\cite{zheng2026}; that claim dates from the April 2026 version of the paper, whose 2022 first posting described a 32-site chain as comparable to the fastest supercomputer.

\textbf{Classical side and evidence.} Tensor-network contraction produced Sycamore-sized samples at comparable XEB values after 24 months~\cite{pan2022}. The boson-sampling experiments are recorded one by one. Mock-up samplers built from low-order marginals came closer than the experiment to the ideal marginals, in total variation distance, for the 2020 100-mode data and for Jiuzhang~2.0, \eld{E34} and \eld{E02} days after those claims~\cite{villalonga2021}. A classically sampleable optical state was not ruled out as an explanation of the 2020 data, while the same tests favoured the quantum model for Jiuzhang~2.0~\cite{martinez2023}. These studies challenge particular validation statistics; they do not sample from the complete ideal distribution. A tensor-network algorithm later outperformed Jiuzhang~2.0~\cite{zhong2021}, Jiuzhang~3.0~\cite{deng2023} and Borealis~\cite{madsen2022} on the statistics used to validate them~\cite{oh2024}, \eld{E36} days after Jiuzhang~3.0 and \eld{E35} days after Borealis. A polynomial-time classical sampler now exists for the distribution used by graph applications of boson sampling; that is a result about that distribution and not about these experiments~\cite{anand2026}. For IBM's 70-qubit circuit, Manabe, Gu and Pan computed the amplitudes of all 2{,}051 published batches of bitstrings in \dcMin{} minutes on 256 GPUs, \eld{E24} days after the claim; the largest intermediate tensor was 256 times smaller than the experiment's authors had estimated~\cite{manabe2026}. The four matches are to a validation benchmark and the two challenges to a diagnostic of one, and each date in Fig.~\ref{fig:timeline} is that of the first response at the validation statistic, which need not be the first cost-competitive reproduction. Amplitudes of published outputs are not fresh samples, and none of these responses shows that a full output distribution can be sampled efficiently.

\textbf{Unresolved.} IBM replaced the circuit in September 2026 with a 64-qubit version of depth 73 with 314 T gates on a ring of qubits; its text states that periodic boundaries remove the contraction path the classical response used~\cite{martiel2026}. We found no classical response to that instance, to the ultracold-atom experiment, to later random-circuit experiments or to the newest boson-sampling generation. Hangleiter concludes that advantage holds for finite-fidelity random-circuit sampling at the level of experimental physics~\cite{hangleiter2026}. The ultracold-atom experiment is validated by Bayesian tests against alternative hypotheses, which do not certify closeness to the ideal distribution.

\textbf{Next test.} A classical response to either revised instance, or sampling at a stated fidelity. An unanswered instance is not an established hardness result, and the three doped Clifford records (original, response, revision) should not be merged into ``refuted'' or ``unrefuted''.

\subsection{Dynamics claims that were answered}
\textbf{Kicked Ising model on 127 qubits (E03).} IBM's Eagle processor ran up to 20 steps of a periodically kicked transverse-field Ising model on its heavy-hexagon lattice and estimated magnetizations by zero-noise extrapolation~\cite{kim2023}. The output is a set of local expectation values. Within 12 days, belief-propagation tensor networks matched to the lattice geometry gave results more accurate than the mitigated data~\cite{tindall2024,aps2024}, and sparse Pauli dynamics reproduced them on one laptop core~\cite{begusic2023,begusic2024}; a dissipative mean-field picture accounts for the experiment's resilience to noise~\cite{dallatorre}. The classical calculations targeted the ideal circuit, so this is reproduction of observables. The explanation is structure: correlations on this lattice are close to tree-like, where belief propagation is exact, and operators grow slowly in the parameter regime studied. A related fault-tolerant target, 100 spins for ten Trotter steps, has a published estimate of 230 logical qubits and $2.4\times10^6$ T gates per execution with repetitions uncosted~\cite{beverland2022}; a converged classical result at that size and depth would answer it.

\textbf{Spin-glass quenches on an annealer (E04).} D-Wave's Advantage2 annealer followed the quench dynamics of Ising spin glasses on two-, three- and infinite-dimensional lattices and reported agreement with the Schr\"odinger equation beyond the reach of the classical methods it tested~\cite{king2025}. A year later, lattice-matched tensor networks~\cite{tindall2025} and time-dependent variational Monte Carlo with Jastrow--Feenberg wavefunctions on up to 128 spins~\cite{mauron2025} reached lower errors than the device in some regimes. Full-state and fourth-order observables were not reproduced, and coverage by lattice, quench time and observable is partial~\cite{kingreply}. The claim is partly answered and contested. Stable long-time variational dynamics in three dimensions at the claimed scales would close it from the classical side; observables beyond classical coverage at matched error would close it from the quantum side.

\subsection{Dynamics claims that stand}
Four claims have no published classical reproduction in their hardest regime. In each, the status is scoped to the methods tested and the criterion stated. It is not a hardness theorem, and disagreement among classical methods shows classical difficulty, not correctness of the quantum result.

\textbf{Out-of-time-order correlators (E06).} An out-of-time-order correlator measures how a local disturbance spreads through a quantum system. Google measured a second-order correlator on 65 qubits of its Willow processor~\cite{google2025otoc}. The primary paper reports 2.1 hours of data collection per circuit and estimates about 3.2 years for tensor-network simulation of one circuit on the Frontier supercomputer, a reported ratio of roughly 13{,}000~\cite{google2025otoc,googleblog2025}. This compares a measured acquisition time with a modelled classical cost. At 40 qubits the measured correlator is checked against exact simulation; at 65 qubits its accuracy is estimated from an empirical error model, and the comparison has not been independently reproduced. We list the two figures and do not plot a ratio. A test by Google co-authors finds the circuits largely incompressible for tensor-network evolution in the Schr\"odinger picture~\cite{bermejo2026}. Open for \eld{E06} days.

\textbf{Floquet magnets on two platforms (E05, E23).} Quantinuum's H2 trapped-ion processor simulated a periodically driven Ising model on a $7\times8$ lattice of 56 qubits; tensor-network, neural-network and operator-truncation methods were tested and none converged in the hardest regime~\cite{haghshenas2025}. Open for \eld{E05} days. IBM's Heron processor, with the QESEM error-mitigation software, measured magnetization on heavy-hexagon ladders of up to 74 qubits and resolved long-lived oscillations~\cite{leviatan2026}. The authors report that tensor networks do not converge and that sparse Pauli-path results depend on truncation, including runs on GPUs and on Fugaku. Their validation has several layers: an unbiased and an extrapolated estimator agree over the time window they share; the deepest and largest results rely on extrapolation; trapped-ion hardware corroborates selected cycles; and the conclusion that oscillations persist in the thermodynamic limit is a finite-size extrapolation. Three questions are therefore separate: whether the trajectory was measured reliably in the tested domain, whether a tested classical method converged there, and whether the extrapolation holds. Open for \eld{E23} days. This is a different task from the 2023 Eagle experiment, and the earlier classical response does not answer it.

\textbf{Operator Loschmidt echo (E07).} IBM and Algorithmiq measured an echo observable on 56 qubits~\cite{barron2026,ibmalgo2026}. Three classical groups produced predictions that disagree with one another. The quantum evidence has two parts of different strength: a heuristic rescaling argument supported by consistency tests across five noise settings on two processors, and a conditional route to error bounds that relies on noise characterization.

\textbf{What would settle these.} From the classical side, a converged reproduction of the stated observable at the stated depth and precision. From the quantum side, validation across the hard regime with the measurement overhead of error mitigation included. Error mitigation estimates a less noisy result from additional noisy runs and a noise model; its bias and its cost both matter. No fault-tolerant compilation exists for any of these observables.

\subsection{Optimization claims that are contested}
\textbf{Annealing with error suppression (E18).} D-Wave hardware with quantum annealing correction, on more than 1{,}300 error-suppressed encoded variables, showed better scaling of time-to-solution than parallel tempering with cluster moves for approximate solutions within 1\% of optimal~\cite{munoz2025,dwave_qac}. A follow-up with the simulated bifurcation machine, a classical dynamical solver, finds comparable or better scaling and concludes that current annealers are unlikely to show a clear advantage on these problem families under its accounting~\cite{pawlowski2026}. The two sides use different runtime conventions and solver portfolios, and that difference is the substance of the dispute.

\textbf{Maximum independent set on Rydberg arrays (E27).} A Harvard experiment on up to 289 atoms reported a superlinear speedup over simulated annealing on the hardest graphs~\cite{ebadi2022}. Andrist et al.\ then solved quasi-planar instances with thousands of nodes to optimality within minutes on commodity hardware and found a relaxed annealer competitive with the quantum algorithm on those instances; the same authors propose harder instances~\cite{andrist2023}.

Both cases teach the lesson of the lattice models in Section~\ref{sec:lattice}: a wider classical portfolio changes an advantage claim. The test that would settle either is a scaling separation from the full portfolio on a defined instance family, with one runtime convention, a stated quality gap and wall time reported. Other proposals fit neither search nor annealing, such as quantum-proposed moves inside a classical Markov chain~\cite{layden2023,orfi2024,marshall2026,hsieh2026}; these need time per statistically independent sample as their metric (class 37).

\subsection{Hybrid workflows that claim no advantage}
\textbf{Iron--sulfur clusters and protein--ligand complexes (E09, E28).} In sample-based quantum diagonalization, a quantum processor samples electronic configurations and a supercomputer diagonalizes the Hamiltonian in the subspace they span. Experiments on IBM Heron processors with Fugaku reached a [4Fe-4S] cluster with 77 qubits, and a closed-loop run used all 152{,}064 Fugaku nodes~\cite{shirakawa2025}. A 2026 study embeds such fragment solvers in two protein--ligand complexes of about 12{,}000 atoms, using up to 94 qubits, \merzCircuits{} circuits and \merzQpu{} processor-hours~\cite{merz2026}. Neither claims an advantage, and both demonstrate integration at scale. Three classical results bound what the quantum part contributes. The energy of the single-layer circuits used can be evaluated classically in polynomial time, reproducing the 77-qubit energy on a laptop in under a minute~\cite{belagali2026}. Sampling from an ideal ground state yields less compact expansions than classical selection: 50{,}551 determinants against 4{,}260 at a 0.1\,Ha error for [2Fe-2S]~\cite{reinholdt2025}. And GPU-accelerated DMRG reaches the [4Fe-4S] active-space size in \feDmrgH{} hours on one node~\cite{legeza2026}, although we have not confirmed that its Hamiltonian is identical to the experiment's. In the protein--ligand study, fragment energies lie 0.10 to 2.3\,mHa above DMRG and classical selected configuration interaction is more accurate on one fragment. Agreement of a binding energy, the error of the fragment solver and the contribution of the quantum sampler are three separate tests; the study reports the second, presents a binding prediction and does not test the third. Further classical results on these subspace methods are reviewed in~\cite{bonilla2026}. What is missing in both is an ablation: the same selection, recovery and diagonalization pipeline fed with classical samples in place of quantum ones.

\textbf{Protein fragments and drug candidates (E15, E16).} Variational quantum folding of short peptides on a lattice produced a structure 1.781\,\AA{} from experiment where the exact optimum of the same lattice model gave 1.879\,\AA~\cite{robert2021,doga2024}. That is not a computational advantage: the device did not optimize the model better than an exact solver, and the model's optimum need not be the most accurate structure. A 2026 study reports better fragment structures than AlphaFold3 under identical post-processing~\cite{zhang2026}; whether classical samplers under the same selection rule match it is untested. A quantum generative prior for KRAS inhibitors gave a modest benchmark gain and two promising synthesized molecules, and its authors state that this does not show superiority over classical priors~\cite{kras2025,uoft2025}. In all three, the quantum contribution has not been isolated.

\subsection{What the device record shows}
Three observations follow. First, every successful response in this set used structured numerical or classical variational methods. No neural network was needed, which is a statement about these responses and not about neural methods (Section~\ref{sec:ai}). Second, a response answers one instance. The doped Clifford and ultracold-atom cases show how a revised experiment restarts the clock, and the Floquet cases show that an answer on one lattice and depth does not carry to another. Third, the standing claims are all coherent dynamics or sampling at depths where known classical representations fail to converge, and each rests on a validation step whose accuracy in that regime is what is in question. The step differs by case: error mitigation or an error model calibrated at smaller sizes for the digital dynamics experiments, error detection with post-selection for the doped Clifford circuits, and statistical tests against alternative hypotheses for the analog sampler. An analog experiment adds a different caution: a Rydberg-atom spin liquid was later simulated variationally at 288 atoms, and the simulated state shows topological order but lacks the entanglement-entropy signature of the resonating-valence-bond state usually associated with it~\cite{semeghini2021,mauron2025tee}; the entropy itself was not measured, although randomized measurements could estimate it~\cite{brydges2019}, so classical simulation there complements the experiment and does not replace it.
